\chapter*{Capstone Project: From Raw Data to Defensible Evidence}
\phantomsection
\addcontentsline{toc}{chapter}{Capstone Project}
\markboth{Capstone Project}{}

Build one reproducible project from raw data to validated conclusions and a reusable model.

\begin{coursechoice}
Use a public dataset not used as a main lab dataset; Heart Failure is allowed. Keep raw files in \texttt{all\_datasets/} and notebooks in \texttt{all\_experiments/}.
\end{coursechoice}

\section*{Required capstone sequence}

\begin{enumerate}
  \item \textbf{Problem.} Specify rows, objective, use moment, population, metric, and error costs.
  \item \textbf{Data.} Record provenance, schema, units, missingness, duplicates, invalid values, identifiers, and cleaning rules.
  \item \textbf{Splits.} Match future use with random, stratified, grouped, or chronological splits; reserve a supervised test set.
  \item \textbf{Baseline.} Evaluate a no-skill or domain rule.
  \item \textbf{Preprocessing.} Fit learned transformations within training folds.
  \item \textbf{Models.} Compare at least three families, including one interpretable model.
  \item \textbf{Development.} Tune, select features, and resample only within validated development pipelines.
  \item \textbf{Diagnostics.} Plot relevant errors, residuals, thresholds, learning curves, features, clusters, or scores.
  \item \textbf{Subgroups.} Locate failures beyond aggregate metrics.
  \item \textbf{Final evidence.} Test once on untouched supervised data, or present complementary unsupervised evidence.
  \item \textbf{Limits.} Distinguish prediction from causation and state unsupported claims.
  \item \textbf{Reproducibility.} Record paths, packages, seeds, notebook order, and configuration.
  \item \textbf{Artifact.} Save the pipeline and model card under \texttt{models/}; demonstrate reloading and prediction.
\end{enumerate}